% ═══════════════════════════════════════════════════════════════
% LERNBLATT: Las preposiciones de lugar
% Spanisch Klasse 8 - Unidad 6, DS 2 (Einzelstunde)
% ═══════════════════════════════════════════════════════════════

\documentclass[11pt, a4paper]{article}

\usepackage[utf8]{inputenc}
\usepackage[T1]{fontenc}
\usepackage[ngerman]{babel}

% --- SCHRIFTEN ---
\usepackage{helvet}
\renewcommand{\familydefault}{\sfdefault}
\usepackage{palatino}
\newcommand{\seriftext}[1]{{\fontfamily{ppl}\selectfont #1}}

\usepackage{microtype}
\usepackage[margin=1.8cm, top=2cm, bottom=1.5cm]{geometry}
\usepackage{enumitem}
\usepackage{tabularx}
\usepackage{booktabs}
\usepackage{array}
\usepackage{multicol}
\usepackage{tikz}
\usetikzlibrary{positioning, shapes.geometric}

\usepackage{fancyhdr}
\usepackage{lastpage}

% ═══════════════════════════════════════════════════════════════
% VARIABLEN
% ═══════════════════════════════════════════════════════════════

\newcommand{\lektion}{06b}
\newcommand{\thema}{Las preposiciones de lugar}
\newcommand{\untertitel}{Die Ortspräpositionen}
\newcommand{\klasse}{8}
\newcommand{\lernziele}{10 Ortspräpositionen · Positionen beschreiben · Stadtviertel beschreiben}

% ═══════════════════════════════════════════════════════════════
% FARBEN
% ═══════════════════════════════════════════════════════════════

\usepackage{xcolor}

\definecolor{colorrojo}{HTML}{C62828}
\definecolor{colorazul}{HTML}{1565C0}
\definecolor{colorverde}{HTML}{2E7D32}
\definecolor{colorgris}{HTML}{757575}
\definecolor{colornegro}{HTML}{222222}

\definecolor{hellgrau}{HTML}{F5F5F5}
\definecolor{rahmen}{HTML}{CCCCCC}
\definecolor{akzent}{HTML}{666666}

% ═══════════════════════════════════════════════════════════════
% BOXEN
% ═══════════════════════════════════════════════════════════════

\usepackage{tcolorbox}
\tcbuselibrary{skins}

\newtcolorbox{lernzielbox}{
    colback=hellgrau, colframe=hellgrau, boxrule=0pt, arc=0pt,
    left=10pt, right=10pt, top=8pt, bottom=8pt,
    before skip=6pt, after skip=8pt
}

\newtcolorbox{grammatikbox}[1][]{
    colback=white, colframe=white, boxrule=0pt, arc=0pt,
    left=10pt, right=6pt, top=6pt, bottom=6pt,
    borderline west={3pt}{0pt}{akzent},
    before skip=4pt, after skip=4pt
}

\newtcolorbox{merkbox}[1][]{
    colback=hellgrau, colframe=hellgrau, boxrule=0pt, arc=0pt,
    left=10pt, right=10pt, top=8pt, bottom=8pt,
    borderline north={2pt}{0pt}{colornegro},
    before skip=6pt, after skip=6pt
}

\newtcolorbox{tippbox}{
    colback=white, colframe=white, boxrule=0pt, arc=0pt,
    left=8pt, right=4pt, top=3pt, bottom=3pt,
    borderline west={2pt}{0pt}{akzent}
}

% ═══════════════════════════════════════════════════════════════
% BEFEHLE
% ═══════════════════════════════════════════════════════════════

\newcommand{\es}[1]{\seriftext{\textbf{#1}}}
\newcommand{\de}[1]{\textit{\textcolor{akzent}{#1}}}
\newcommand{\gram}[1]{\textsc{#1}}
\newcommand{\abmark}[1]{{\footnotesize\textcolor{akzent}{$\rightarrow$ AB #1}}}
\newcommand{\abschnitt}[2]{\noindent\textbf{\large #1. #2}}

\setlength{\parindent}{0pt}
\setlength{\parskip}{5pt}

% ═══════════════════════════════════════════════════════════════
% KOPF- UND FUSSZEILE
% ═══════════════════════════════════════════════════════════════

\pagestyle{fancy}
\fancyhf{}
\renewcommand{\headrulewidth}{0.5pt}
\renewcommand{\footrulewidth}{0pt}
\lhead{\footnotesize\textsc{Spanisch} · Klasse \klasse}
\rhead{\footnotesize\thema}
\cfoot{\footnotesize\thepage\,/\,\pageref{LastPage}}

% ═══════════════════════════════════════════════════════════════
% DOKUMENT
% ═══════════════════════════════════════════════════════════════

\begin{document}

% --- TITEL ---
\begin{center}
    {\LARGE\bfseries \thema}\\[3pt]
    {\small \untertitel}
\end{center}

\vspace{4pt}

% --- EINLEITUNG ---
Wo ist die Bäckerei? Neben dem Park oder gegenüber vom Kino?
Mit den Ortspräpositionen kannst du beschreiben, wo sich Dinge und Orte befinden.

\vspace{2pt}

% --- LERNZIELE ---
\begin{lernzielbox}
\textbf{Das lernst du:}\quad \lernziele
\end{lernzielbox}

% ═══════════════════════════════════════════════════════════════
% HAUPTTABELLE: PRÄPOSITIONEN
% ═══════════════════════════════════════════════════════════════

\begin{merkbox}
\textbf{Die 10 wichtigsten Ortspräpositionen}

\vspace{6pt}

\begin{center}
\renewcommand{\arraystretch}{1.4}
\begin{tabular}{|l|l|l|}
\hline
\textbf{Español} & \textbf{Deutsch} & \textbf{Beispiel} \\
\hline
\es{enfrente de} & gegenüber von & El cine está enfrente de la plaza. \\
\hline
\es{encima de} & über, auf & El gato está encima de la mesa. \\
\hline
\es{al lado de} & neben & La farmacia está al lado del parque. \\
\hline
\es{a la izquierda de} & links von & El banco está a la izquierda del teatro. \\
\hline
\es{en} & in, an, auf & Estoy en casa. / Estoy en la calle. \\
\hline
\es{a la derecha de} & rechts von & La iglesia está a la derecha del cine. \\
\hline
\es{entre} & zwischen & El parque está entre el cine y la iglesia. \\
\hline
\es{detrás de} & hinter & El coche está detrás de la casa. \\
\hline
\es{debajo de} & unter & El perro está debajo de la silla. \\
\hline
\es{delante de} & vor & El árbol está delante de la ventana. \\
\hline
\end{tabular}
\end{center}
\end{merkbox}

\vspace{4pt}

\abmark{Tabelle übertragen}

\vspace{6pt}

% ═══════════════════════════════════════════════════════════════
% ZWEI SPALTEN
% ═══════════════════════════════════════════════════════════════

\begin{multicols}{2}

% ---------------------------------------------------------------
% VISUALISIERUNG
% ---------------------------------------------------------------
\abschnitt{1}{Visualisierung}

\vspace{4pt}

\begin{center}
\begin{tikzpicture}[scale=0.9]
    % Tisch
    \draw[thick] (-1.5,-0.5) rectangle (1.5,0);
    \node at (0,-0.25) {\small mesa};

    % Beine
    \draw[thick] (-1.3,-0.5) -- (-1.3,-1.5);
    \draw[thick] (1.3,-0.5) -- (1.3,-1.5);

    % Ball auf dem Tisch
    \fill[colorazul] (0,0.4) circle (0.3);
    \node[right] at (0.5,0.4) {\footnotesize encima de};

    % Ball unter dem Tisch
    \fill[colorrojo] (0,-1) circle (0.3);
    \node[right] at (0.5,-1) {\footnotesize debajo de};

    % Ball links
    \fill[colorverde] (-2.2,-0.25) circle (0.3);
    \node[left] at (-2.5,-0.25) {\footnotesize a la izquierda};

    % Ball rechts
    \fill[colorverde] (2.2,-0.25) circle (0.3);
    \node[right] at (2.5,-0.25) {\footnotesize a la derecha};
\end{tikzpicture}
\end{center}

\vspace{6pt}

\begin{tippbox}
\footnotesize\textit{Merke:} \es{de + el = del}\\
"al lado \textbf{del} parque" (nicht: *de el parque)
\end{tippbox}

\columnbreak

% ---------------------------------------------------------------
% REDEMITTEL
% ---------------------------------------------------------------
\abschnitt{2}{Redemittel}

\vspace{4pt}

\begin{grammatikbox}
\textbf{Die Struktur:}\\[2pt]
\es{El/La [Ort] está [Präposition] [Bezugspunkt].}
\end{grammatikbox}

\vspace{4pt}

\textbf{Beispiele:}

\begin{itemize}[leftmargin=1.2em, itemsep=2pt]
    \item \es{La biblioteca está enfrente del parque.}\\
          \de{Die Bibliothek ist gegenüber vom Park.}
    \item \es{El supermercado está al lado de la farmacia.}\\
          \de{Der Supermarkt ist neben der Apotheke.}
    \item \es{La iglesia está entre el cine y el teatro.}\\
          \de{Die Kirche ist zwischen dem Kino und dem Theater.}
\end{itemize}

\vspace{4pt}

\abmark{Aufgabe 1 + 2}

\end{multicols}

\vspace{4pt}

% ═══════════════════════════════════════════════════════════════
% BUCHÜBUNGEN
% ═══════════════════════════════════════════════════════════════

\begin{grammatikbox}
\textbf{Buchübungen (Qué pasa 1)}

\vspace{4pt}

\textbf{Übung 2B (S. 106) -- Copito:}\\
Schau dir die 10 Bilder von Copito an. Beschreibe für jedes Bild, wo der Hund ist.\\
\textit{Beispiel:} Bild 1 $\rightarrow$ "Copito está debajo de la mesa."

\vspace{6pt}

\textbf{Übung 4 (S. 107) -- Dos barrios:}\\
1. Hört das Audio. Welches Viertel wird beschrieben -- A oder B?\\
2. Beschreibt danach das ANDERE Viertel mit mindestens 5 Sätzen.\\
\textit{Redemittel:} El/La \_\_\_ está enfrente de... / al lado de... / entre... y...
\end{grammatikbox}

\vspace{6pt}

% ═══════════════════════════════════════════════════════════════
% KULTURELLES
% ═══════════════════════════════════════════════════════════════

\begin{tippbox}
\textbf{¿Sabías que...?} In spanischen Städten gibt es oft eine \es{plaza mayor} (Hauptplatz) im Zentrum. Alle wichtigen Gebäude (Rathaus, Kirche, Cafés) sind um diesen Platz herum angeordnet. Man sagt: "La iglesia está al lado del ayuntamiento" (Die Kirche ist neben dem Rathaus).
\end{tippbox}

\end{document}
